\documentclass[11pt,a4paper]{report}
\usepackage[utf8]{inputenc}
\usepackage[portuguese]{babel}
\usepackage[margin=2.5cm]{geometry}
\usepackage{graphicx}
\usepackage{booktabs}
\usepackage{amsmath}
\usepackage{hyperref}
\usepackage{listings}
\usepackage{xcolor}
\usepackage{titlesec}
\usepackage{float}
\usepackage{indentfirst}

\titleformat{\chapter}
  {\normalfont\huge\bfseries}{}{0em}{}
\titlespacing*{\chapter}{0pt}{0pt}{20pt}

\setlength{\parskip}{0.6em} 
\setlength{\parindent}{15pt}

\begin{document}

\begin{titlepage}
    \centering
    
    {\Large {Universidade de Aveiro}\par}


    \vspace{0.5cm}

    {\Large \emph{Departamento de Eletrónica, Telecomunicações e Informática}\par}

    \vfill
    
    {\Huge \textbf{Lab 1}\par}
    \vspace{0.5cm}
    {\Large Informação e Codificação\par}

    \vspace{1cm}
    \includegraphics[width=0.2\textwidth]{ua_logo.png}\par
    
    \vspace{3cm}
    
    André Reis (119814)\par
    Guilherme Sousa (119080)\par
    João Ferreira (Número / ID)\par

    \vfill

    {\large Data\par}
\end{titlepage}

\begin{abstract}
Resumo.
\end{abstract}


\tableofcontents
\newpage

\section{Introdução}
Introdução


\chapter{Parte I}

\section{Exercício 1}

\section{Exercício 2: Quantização Uniforme (wav\_quant)}

O objetivo deste exercício consistiu no desenvolvimento de um programa, denominado \emph{wav\_quant}, capaz de reduzir o número de bits efetivos por amostra de um ficheiro de áudio no formato WAV, através de quantização escalar uniforme.

\subsection{Implementação}

Tal como no exercício 4, optámos por usar Python com a biblioteca \emph{numpy}. O módulo \emph{wave} lê e escreve no contentor WAV, e o numpy permite aplicar a quantização a todas as amostras de uma só vez, através de operações vetoriais.

A quantização opera diretamente sobre os valores inteiros, pelo que os dados são convertidos para inteiros com sinal de 64 bits e organizados numa matriz $(N, C)$ com $N$ fotogramas e $C$ canais. As amostras de 8 bits, que no formato WAV são sem sinal, são centradas em zero subtraindo 128; as de 24 bits, para as quais não existe tipo nativo no numpy, são reconstruídas a partir dos 3 bytes com deslocamentos de bits, seguida de extensão de sinal.

A entrada/saída (read\_wav, write\_wav) está separada da quantização. O programa suporta os canais mono, estéreo e multicanal e amostras de 8, 16, 24 e 32 bits, preservando a largura original na escrita. Assim, o ficheiro de saída tem o mesmo tamanho que o de entrada: apenas o número de níveis distintos utilizados é reduzido.

\subsection{Quantização}

Sendo $N_b$ a resolução original e $B$ o número de bits pretendido, descartam-se os $N_b - B$ bits menos significativos. O passo de quantização é
\[
\Delta = 2^{\,N_b - B},
\]
e os índices de quantização válidos pertencem ao intervalo $[-2^{B-1},\, 2^{B-1}-1]$. Se $B = N_b$, o sinal é devolvido sem alterações. Foram implementados três quantizadores:

\begin{enumerate}
    \item \textbf{Arredondamento (round):}
    Cada amostra é aproximada ao nível de quantização mais próximo:
    \[
    q[n] = \left\lfloor \frac{x[n]}{\Delta} + \frac{1}{2} \right\rfloor, \qquad y[n] = q[n]\,\Delta,
    \]
    com $q[n]$ limitado ao intervalo de índices válidos, para não ultrapassar o último nível. O erro está contido em $[-\Delta/2, \Delta/2]$. É o modo por omissão e, em geral, o que produz menor erro.

    \item \textbf{Truncatura (truncate):}
    Os bits menos significativos são simplesmente descartados, o que equivale a arredondar sempre para baixo:
    \[
    y[n] = \left\lfloor \frac{x[n]}{\Delta} \right\rfloor \Delta.
    \]
    Na implementação, isto corresponde a um deslocamento aritmético para a direita seguido de um deslocamento para a esquerda de $N_b - B$ bits. O erro está em $(-\Delta, 0]$, o que introduz um desvio médio (componente DC) de aproximadamente $-\Delta/2$.

    \item \textbf{Mid-rise (midrise):}
    Trunca e soma meio passo, colocando o valor reconstruído no centro de cada intervalo:
    \[
    y[n] = \left\lfloor \frac{x[n]}{\Delta} \right\rfloor \Delta + \frac{\Delta}{2}.
    \]
    Elimina o desvio da truncatura e, como é característico dos quantizadores \emph{mid-rise}, não possui nível de reconstrução em zero.
\end{enumerate}

Em todos os casos, o resultado é limitado ao intervalo representável com $N_b$ bits, evitando overflow na escrita. A Tabela~\ref{tab:quant_exemplo} ilustra os três modos para uma amostra de 16 bits reduzida a 8 bits ($\Delta = 256$).

\begin{table}[H]
\centering
\begin{tabular}{lcc}
\toprule
Modo & Cálculo & $y$ \\
\midrule
round & $\lfloor 1000/256 + 0.5 \rfloor \cdot 256 = 4 \cdot 256$ & 1024 \\
truncate & $\lfloor 1000/256 \rfloor \cdot 256 = 3 \cdot 256$ & 768 \\
midrise & $768 + 128$ & 896 \\
\bottomrule
\end{tabular}
\caption{Quantização da amostra $x = 1000$ de 16 para 8 bits}
\label{tab:quant_exemplo}
\end{table}

Como regra geral, cada bit removido degrada a relação sinal-ruído de quantização em cerca de 6~dB. Abaixo de aproximadamente 8 bits o ruído torna-se claramente audível e, com 1 a 3 bits, o sinal fica fortemente distorcido.

\subsection{Utilização}

\texttt{pip install numpy}

\texttt{python3 wav\_quant.py input.wav output.wav bits [modo]}

O argumento \texttt{bits} deve estar entre 1 e a resolução do ficheiro de entrada. O modo é opcional e, por omissão, é \texttt{round}.

\begin{table}[H]
\centering
\label{tab:modos_quant}
\begin{tabular}{lc}
\toprule
Quantizador & Modo \\
\midrule
Arredondamento & \texttt{round} \\
Truncatura & \texttt{truncate} \\
Mid-rise & \texttt{midrise} \\
\bottomrule
\end{tabular}
\caption{Modos de quantização disponíveis}
\end{table}

\section{Exercício 3}

\section{Exercício 4: Implementação de Efeitos de Áudio (wav\_effects)}

O objetivo deste exercício consistiu no desenvolvimento de um programa, denominado \emph{wav\_effects}, capaz de aplicar transformações acústicas a ficheiros de áudio PCM no formato WAV. 

\subsection{Implementação}

Optámos por usar Python com a biblioteca \emph{numpy}. O módulo \emph{wave} da biblioteca lê e escreve no contentor WAV, e o numpy permite expressar os efeitos como operações sobre arrays, com um código curto e legível e um desempenho.

Os dados são convertidos para vírgula flutuante de 32 bits no intervalo $[-1,1]$, organizados numa matriz $(N, C)$ com $N$ fotogramas e $C$ canais. Esta representação torna os efeitos independentes da largura de amostra original e permite somar sinais sem overflow de inteiros.

A entrada/saída (read\_wav, write\_wav) está separada dos efeitos.

 O programa suporta os canais mono, estéreo e multicanal e amostras PCM de 8, 16, 24 e 32 bits, preservando a largura original na escrita.

 \subsection{Efeitos}

O sistema divide-se em funções modulares autónomas para cada efeito:

\begin{enumerate}
    \item \textbf{Eco Simples:}
    Soma ao sinal original uma cópia atrasada e atenuada:
    \[
    y[n] = x[n] + a\,x[n-D], \qquad D = \lfloor \tau f_s \rfloor,
    \]
    com $\tau = 0.3$~s e $a = 0.7$. O buffer de saída tem $N + D$ amostras, para acomodar a cauda do eco, pelo que o ficheiro fica $0.3$~s mais longo.

    \item \textbf{Ecos Múltiplos:}
    Generaliza o eco simples com $K$ repetições cuja amplitude decai geometricamente:
    \[
    y[n] = \sum_{k=0}^{K} a^{k}\,x[n-kD],
    \]
    com $\tau = 0.25$~s, $a = 0.7$ e $K = 6$, ou seja, o sinal original seguido de seis ecos. A duração aumenta $K\tau = 1.5$~s. Como as repetições se sobrepõem quando o sinal é mais longo do que $D$, as amplitudes somam-se.

    \item \textbf{Modulação de Amplitude:}
     Multiplica o sinal por um oscilador de baixa frequência (LFO) sinusoidal, produzindo um efeito de tremolo:
    \[
    y[n] = x[n]\,m[n], \qquad m[n] = 1 - \frac{p}{2} + \frac{p}{2}\sin\!\left(2\pi f_m \frac{n}{f_s}\right),
    \]
    com profundidade $p = 0.8$ e $f_m = 4$~Hz. O ganho varia entre $1-p = 0.2$ e $1$. A duração mantém-se.
    
    \item \textbf{Atraso Variável no Tempo:}
    Soma ao sinal original uma cópia cujo atraso varia continuamente:
    \[
    y[n] = x[n] + 0.6\,x[n - d[n]], \qquad d[n] = D_0 + A\sin\!\left(2\pi f_m \frac{n}{f_s}\right),
    \]
    com $D_0 = 15$~ms, $A = 10$~ms e $f_m = 0.5$~Hz (o atraso oscila entre 5 e 25~ms). Este processo é conhecido como \emph{flanger}: a combinação do sinal com uma cópia de atraso curto é um filtro em pente, e a variação do atraso desloca os cortes espectrais ao longo do tempo. Como $s = n - d[n]$ é em geral fracionário, o valor é obtido por interpolação linear:
    \[
    x[s] \approx (1-\phi)\,x[\lfloor s \rfloor] + \phi\,x[\lfloor s \rfloor + 1], \qquad \phi = s - \lfloor s \rfloor.
    \]
    Arredondar $s$ para o inteiro mais próximo produziria saltos bruscos no atraso, audíveis como cliques e \emph{zipper noise}. Nas amostras iniciais, onde $s<0$, apenas se mantém $x[n]$. A duração do ficheiro não se altera.
\end{enumerate}

\subsection{Utilização}

\texttt{pip install numpy}

\texttt{python3 wav\_effects.py input.wav output.wav sigla}

\begin{table}[H]
\centering
\label{tab:efeitos}
\begin{tabular}{lc}
\toprule
Efeito & Sigla \\
\midrule
Eco Simples & \texttt{se} \\
Ecos Múltiplos & \texttt{me} \\
Modulação de Amplitude & \texttt{am} \\
Atraso Variável no Tempo & \texttt{vd} \\
\bottomrule
\end{tabular}
\caption{Efeitos de áudio disponíveis}
\end{table}


\chapter{Parte II}

\section{Exercício 5}

\section{Exercício 6}


\chapter{Parte III}

\section{Exercício 7}

\section{Exercício 8}


\chapter{Conclusão}
Conclusão.


\end{document}\documentclass[11pt,a4paper]{report}
\usepackage[utf8]{inputenc}
\usepackage[portuguese]{babel}
\usepackage[margin=2.5cm]{geometry}
\usepackage{graphicx}
\usepackage{booktabs}
\usepackage{amsmath}
\usepackage{hyperref}
\usepackage{listings}
\usepackage{xcolor}
\usepackage{titlesec}
\usepackage{float}
\usepackage{indentfirst}

\titleformat{\chapter}
  {\normalfont\huge\bfseries}{}{0em}{}
\titlespacing*{\chapter}{0pt}{0pt}{20pt}

\setlength{\parskip}{0.6em}
\setlength{\parindent}{15pt}

\begin{document}

\begin{titlepage}
    \centering

    {\Large {Universidade de Aveiro}\par}


    \vspace{0.5cm}

    {\Large \emph{Departamento de Eletrónica, Telecomunicações e Informática}\par}

    \vfill

    {\Huge \textbf{Lab 1}\par}
    \vspace{0.5cm}
    {\Large Informação e Codificação\par}

    \vspace{1cm}
    \includegraphics[width=0.2\textwidth]{ua_logo.png}\par

    \vspace{3cm}

    André Reis (119814)\par
    Guilherme Sousa (119080)\par
    João Ferreira (Número / ID)\par

    \vfill

    {\large Data\par}
\end{titlepage}

\begin{abstract}
Resumo.
\end{abstract}


\tableofcontents
\newpage

\section{Introdução}
Introdução


\chapter{Parte I}

\section{Exercício 1}

\section{Exercício 2: Quantização Uniforme (wav\_quant)}

O objetivo deste exercício consistiu no desenvolvimento de um programa, denominado \emph{wav\_quant}, capaz de reduzir o número de bits efetivos por amostra de um ficheiro de áudio no formato WAV, através de quantização escalar uniforme.

\subsection{Implementação}

Tal como no exercício 4, optámos por usar Python com a biblioteca \emph{numpy}. O módulo \emph{wave} lê e escreve no contentor WAV, e o numpy permite aplicar a quantização a todas as amostras de uma só vez, através de operações vetoriais.

A quantização opera diretamente sobre os valores inteiros, pelo que os dados são convertidos para inteiros com sinal de 64 bits e organizados numa matriz $(N, C)$ com $N$ fotogramas e $C$ canais. As amostras de 8 bits, que no formato WAV são sem sinal, são centradas em zero subtraindo 128; as de 24 bits, para as quais não existe tipo nativo no numpy, são reconstruídas a partir dos 3 bytes com deslocamentos de bits, seguida de extensão de sinal.

A entrada/saída (read\_wav, write\_wav) está separada da quantização. O programa suporta os canais mono, estéreo e multicanal e amostras de 8, 16, 24 e 32 bits, preservando a largura original na escrita. Assim, o ficheiro de saída tem o mesmo tamanho que o de entrada: apenas o número de níveis distintos utilizados é reduzido.

\subsection{Quantização}

Sendo $N_b$ a resolução original e $B$ o número de bits pretendido, descartam-se os $N_b - B$ bits menos significativos. O passo de quantização é
\[
\Delta = 2^{\,N_b - B},
\]
e os índices de quantização válidos pertencem ao intervalo $[-2^{B-1},\, 2^{B-1}-1]$. Se $B = N_b$, o sinal é devolvido sem alterações. Foram implementados três quantizadores:

\begin{enumerate}
    \item \textbf{Arredondamento (round):}
    Cada amostra é aproximada ao nível de quantização mais próximo:
    \[
    q[n] = \left\lfloor \frac{x[n]}{\Delta} + \frac{1}{2} \right\rfloor, \qquad y[n] = q[n]\,\Delta,
    \]
    com $q[n]$ limitado ao intervalo de índices válidos, para não ultrapassar o último nível. O erro está contido em $[-\Delta/2, \Delta/2]$. É o modo por omissão e, em geral, o que produz menor erro.

    \item \textbf{Truncatura (truncate):}
    Os bits menos significativos são simplesmente descartados, o que equivale a arredondar sempre para baixo:
    \[
    y[n] = \left\lfloor \frac{x[n]}{\Delta} \right\rfloor \Delta.
    \]
    Na implementação, isto corresponde a um deslocamento aritmético para a direita seguido de um deslocamento para a esquerda de $N_b - B$ bits. O erro está em $(-\Delta, 0]$, o que introduz um desvio médio (componente DC) de aproximadamente $-\Delta/2$.

    \item \textbf{Mid-rise (midrise):}
    Trunca e soma meio passo, colocando o valor reconstruído no centro de cada intervalo:
    \[
    y[n] = \left\lfloor \frac{x[n]}{\Delta} \right\rfloor \Delta + \frac{\Delta}{2}.
    \]
    Elimina o desvio da truncatura e, como é característico dos quantizadores \emph{mid-rise}, não possui nível de reconstrução em zero.
\end{enumerate}

Em todos os casos, o resultado é limitado ao intervalo representável com $N_b$ bits, evitando overflow na escrita. A Tabela~\ref{tab:quant_exemplo} ilustra os três modos para uma amostra de 16 bits reduzida a 8 bits ($\Delta = 256$).

\begin{table}[H]
\centering
\begin{tabular}{lcc}
\toprule
Modo & Cálculo & $y$ \\
\midrule
round & $\lfloor 1000/256 + 0.5 \rfloor \cdot 256 = 4 \cdot 256$ & 1024 \\
truncate & $\lfloor 1000/256 \rfloor \cdot 256 = 3 \cdot 256$ & 768 \\
midrise & $768 + 128$ & 896 \\
\bottomrule
\end{tabular}
\caption{Quantização da amostra $x = 1000$ de 16 para 8 bits}
\label{tab:quant_exemplo}
\end{table}

Como regra geral, cada bit removido degrada a relação sinal-ruído de quantização em cerca de 6~dB. Abaixo de aproximadamente 8 bits o ruído torna-se claramente audível e, com 1 a 3 bits, o sinal fica fortemente distorcido.

\subsection{Utilização}

\texttt{pip install numpy}

\texttt{python3 wav\_quant.py input.wav output.wav bits [modo]}

O argumento \texttt{bits} deve estar entre 1 e a resolução do ficheiro de entrada. O modo é opcional e, por omissão, é \texttt{round}.

\begin{table}[H]
\centering
\label{tab:modos_quant}
\begin{tabular}{lc}
\toprule
Quantizador & Modo \\
\midrule
Arredondamento & \texttt{round} \\
Truncatura & \texttt{truncate} \\
Mid-rise & \texttt{midrise} \\
\bottomrule
\end{tabular}
\caption{Modos de quantização disponíveis}
\end{table}

\section{Exercício 3}

\section{Exercício 4: Implementação de Efeitos de Áudio (wav\_effects)}

O objetivo deste exercício consistiu no desenvolvimento de um programa, denominado \emph{wav\_effects}, capaz de aplicar transformações acústicas a ficheiros de áudio PCM no formato WAV.

\subsection{Implementação}

Optámos por usar Python com a biblioteca \emph{numpy}. O módulo \emph{wave} da biblioteca lê e escreve no contentor WAV, e o numpy permite expressar os efeitos como operações sobre arrays, com um código curto e legível e um desempenho.

Os dados são convertidos para vírgula flutuante de 32 bits no intervalo $[-1,1]$, organizados numa matriz $(N, C)$ com $N$ fotogramas e $C$ canais. Esta representação torna os efeitos independentes da largura de amostra original e permite somar sinais sem overflow de inteiros.

A entrada/saída (read\_wav, write\_wav) está separada dos efeitos.

 O programa suporta os canais mono, estéreo e multicanal e amostras PCM de 8, 16, 24 e 32 bits, preservando a largura original na escrita.

 \subsection{Efeitos}

O sistema divide-se em funções modulares autónomas para cada efeito:

\begin{enumerate}
    \item \textbf{Eco Simples:}
    Soma ao sinal original uma cópia atrasada e atenuada:
    \[
    y[n] = x[n] + a\,x[n-D], \qquad D = \lfloor \tau f_s \rfloor,
    \]
    com $\tau = 0.3$~s e $a = 0.7$. O buffer de saída tem $N + D$ amostras, para acomodar a cauda do eco, pelo que o ficheiro fica $0.3$~s mais longo.

    \item \textbf{Ecos Múltiplos:}
    Generaliza o eco simples com $K$ repetições cuja amplitude decai geometricamente:
    \[
    y[n] = \sum_{k=0}^{K} a^{k}\,x[n-kD],
    \]
    com $\tau = 0.25$~s, $a = 0.7$ e $K = 6$, ou seja, o sinal original seguido de seis ecos. A duração aumenta $K\tau = 1.5$~s. Como as repetições se sobrepõem quando o sinal é mais longo do que $D$, as amplitudes somam-se.

    \item \textbf{Modulação de Amplitude:}
     Multiplica o sinal por um oscilador de baixa frequência (LFO) sinusoidal, produzindo um efeito de tremolo:
    \[
    y[n] = x[n]\,m[n], \qquad m[n] = 1 - \frac{p}{2} + \frac{p}{2}\sin\!\left(2\pi f_m \frac{n}{f_s}\right),
    \]
    com profundidade $p = 0.8$ e $f_m = 4$~Hz. O ganho varia entre $1-p = 0.2$ e $1$. A duração mantém-se.

    \item \textbf{Atraso Variável no Tempo:}
    Soma ao sinal original uma cópia cujo atraso varia continuamente:
    \[
    y[n] = x[n] + 0.6\,x[n - d[n]], \qquad d[n] = D_0 + A\sin\!\left(2\pi f_m \frac{n}{f_s}\right),
    \]
    com $D_0 = 15$~ms, $A = 10$~ms e $f_m = 0.5$~Hz (o atraso oscila entre 5 e 25~ms). Este processo é conhecido como \emph{flanger}: a combinação do sinal com uma cópia de atraso curto é um filtro em pente, e a variação do atraso desloca os cortes espectrais ao longo do tempo. Como $s = n - d[n]$ é em geral fracionário, o valor é obtido por interpolação linear:
    \[
    x[s] \approx (1-\phi)\,x[\lfloor s \rfloor] + \phi\,x[\lfloor s \rfloor + 1], \qquad \phi = s - \lfloor s \rfloor.
    \]
    Arredondar $s$ para o inteiro mais próximo produziria saltos bruscos no atraso, audíveis como cliques e \emph{zipper noise}. Nas amostras iniciais, onde $s<0$, apenas se mantém $x[n]$. A duração do ficheiro não se altera.
\end{enumerate}

\subsection{Utilização}

\texttt{pip install numpy}

\texttt{python3 wav\_effects.py input.wav output.wav sigla}

\begin{table}[H]
\centering
\label{tab:efeitos}
\begin{tabular}{lc}
\toprule
Efeito & Sigla \\
\midrule
Eco Simples & \texttt{se} \\
Ecos Múltiplos & \texttt{me} \\
Modulação de Amplitude & \texttt{am} \\
Atraso Variável no Tempo & \texttt{vd} \\
\bottomrule
\end{tabular}
\caption{Efeitos de áudio disponíveis}
\end{table}


\chapter{Parte II}

\section{Exercício 5}

\section{Exercício 6}


\chapter{Parte III}

\section{Exercício 7}

\section{Exercício 8}


\chapter{Conclusão}
Conclusão.


\end{document}